\chapter{Heuristics That Predict}
\label{ch:heuristics}

\begin{center}\itshape
A heuristic is a proof that is wrong in an interesting way.\\
It tells you what is true before anyone can show it,\\
and sometimes it tells you where to look.
\end{center}

Here is a list of identities. Read it slowly:
\[
\begin{aligned}
3&=1^3+1^3+1^3=4^3+4^3+(-5)^3,\\
30&=2\,220\,422\,932^3+(-2\,218\,888\,517)^3+(-283\,059\,965)^3,\\
33&=8\,866\,128\,975\,287\,528^3+(-8\,778\,405\,442\,862\,239)^3+(-2\,736\,111\,468\,807\,040)^3,\\
42&=(-80\,538\,738\,812\,075\,974)^3+80\,435\,758\,145\,817\,515^3+12\,602\,123\,297\,335\,631^3.
\end{aligned}
\]
In the 1950s Louis Mordell asked whether $3$ has any representation as a sum of three cubes other than the two on the first line. The
representation of $30$ was found in 1999, after decades of searching. In 2015 the YouTube channel Numberphile made a video about $33$, the smallest number
for which no representation was known; Andrew Booker of the University of Bristol watched it, designed a new algorithm, and in March 2019, after three weeks on a
university cluster, found the third line. Six months later, Booker and Andrew Sutherland of MIT, with half a million volunteer computers of the
Charity Engine grid, found $42$ (the last case below $100$), and then a third representation of $3$:
\[
3=569\,936\,821\,221\,962\,380\,720^3+(-569\,936\,821\,113\,563\,493\,509)^3+(-472\,715\,493\,453\,327\,032)^3.
\]
Every one of these identities can be checked by hand in an afternoon, and by computer in a microsecond.

Why did anyone believe these solutions exist? Why did they expect them to be so large? And why did they search where they searched? The
answer is a heuristic, due to Roger Heath-Brown, that predicts not only \emph{whether} there are solutions but \emph{how many} there are up to a given size.
This chapter derives the heuristic, proves the parts of it that can be proved, and proves the trick that made Booker's search possible.

\section{What can be proved}

\begin{proposition}\label{prop:mod9}
If $n\equiv\pm4\pmod9$, then $n$ is not a sum of three cubes of integers.
\end{proposition}

\begin{proof}
If $3\mid x$ then $x^3\equiv0\pmod{27}$. Otherwise $x\equiv\pm1,\pm2,\pm4\pmod9$, and $x^3\equiv\pm1,\pm8,\pm64$, all $\equiv\pm1\pmod9$. So cubes are $0,\pm1$ modulo $9$, and a sum of three of them is in
$\{0,\pm1,\pm2,\pm3\}$ modulo $9$, which misses $\pm4$.
\end{proof}

For two numbers, infinitely many representations are known, given by polynomial identities.

\begin{proposition}[Verebrusov, 1908; Mahler, 1936]\label{prop:families}
For every integer $t$,
\[
(1+6t^3)^3+(1-6t^3)^3+(-6t^2)^3=2,\qquad(9t^4)^3+(3t-9t^4)^3+(1-9t^3)^3=1.
\]
\end{proposition}

\begin{proof}
For the first, $(1+u)^3+(1-u)^3=2+6u^2$; with $u=6t^3$ this is $2+216t^6=2+(6t^2)^3$. For the second, let $a=9t^4$ and $b=3t-9t^4$, so $a+b=3t$ and
$a^3+b^3=(a+b)^3-3ab(a+b)=27t^3-9t\cdot9t^4(3t-9t^4)=27t^3-243t^6+729t^9$; and $(1-9t^3)^3=1-27t^3+243t^6-729t^9$. The sum is $1$.
\end{proof}

No such family is known for any other $n$ that is not a cube (a cube $n=m^3$ has the family $m^3+t^3+(-t)^3$), and for $3$ it is not even known whether there are infinitely many
solutions. So the conjecture that every $n\not\equiv\pm4\pmod9$ is a sum of three cubes is open; it is believed, because of the heuristic.

\section{The heuristic}

Count solutions in a box. Let $N_n(B)$ be the number of $(x,y,z)\in\Z^3$ with $x^3+y^3+z^3=n$ and $\max(|x|,|y|,|z|)\le B$. The heuristic has two factors, one from the real numbers and one from the
primes; we derive each, and say exactly which step is not a proof.

\subsection*{The real density}

Pretend that the value $F=x^3+y^3+z^3$ at a random integer point near a point of $\R^3$ is ``spread out'' like a continuous variable. The expected number of integer points with $F=n$ is then the
volume of the region where $|F-n|<\frac12$. For a dyadic shell $B<\max(|x|,|y|,|z|)\le2B$, substitute $x=Bu$, $y=Bv$, $z=Bw$: the region becomes $\{1<\max(|u|,|v|,|w|)\le2,\ |u^3+v^3+w^3-n/B^3|<\frac1{2B^3}\}$, of
volume about $\frac1{B^3}A$, where $A$ is the ``area'' of the surface $u^3+v^3+w^3=0$ in the shell, measured with the density $\frac1{|\nabla F|}$; multiplied by $B^3$ (the Jacobian), the shell
contributes a constant $A$, \emph{independent of $B$}. There are $\log_2B$ dyadic shells up to $B$, so the real part of the heuristic predicts
\[
N_n(B)\approx\sigma_\infty\log B,
\]
with an explicit constant $\sigma_\infty$; following Heath-Brown, Booker and Sutherland find $\sigma_\infty=4\int_1^\infty(t^3+1)^{-2/3}\dd t=\frac{2\,\Gamma(\frac13)^2}{3\,\Gamma(\frac23)}\approx3.533$ (Exercise~\ref{ex:sigmainf} explains where the integral comes from). The logarithm is the
signature of the problem: three variables of degree three is exactly the critical case where each scale contributes the same amount (Exercise~\ref{ex:dimcount}).

\subsection*{The local densities}

The integers are not random: congruences can forbid solutions, as Proposition~\ref{prop:mod9} shows, or favour them. For a prime $p$, the \emph{local density} measures how much more (or less) often
$F\equiv n$ happens modulo powers of $p$ than for a random value:
\[
\sigma_p(n)=\lim_{k\to\infty}\frac{\#\{(x,y,z)\bmod p^k:\ x^3+y^3+z^3\equiv n\pmod{p^k}\}}{p^{2k}}.
\]
(For a random value, the count would be $p^{3k}\cdot p^{-k}=p^{2k}$.)

\begin{proposition}\label{prop:localdens}
\begin{enumerate}[label=(\alph*)]
\item If $p\equiv2\pmod3$ and $p\nmid n$, then $\sigma_p(n)=1$.
\item If $n\equiv\pm4\pmod9$, then $\sigma_3(n)=0$.
\end{enumerate}
\end{proposition}

\begin{proof}
(a) \emph{Modulo $p$.} The group $\mathbb F_p^\times$ has order $p-1$, prime to $3$, so $x\mapsto x^3$ is injective on it, and since $0^3=0$ it is a bijection of $\mathbb F_p$. Hence $(x,y,z)\mapsto(x^3,y^3,z^3)$ permutes $\mathbb F_p^3$, and
the number of solutions of $x^3+y^3+z^3=n$ in $\mathbb F_p$ equals the number of $(a,b,c)$ with $a+b+c=n$, which is $p^2$.

\emph{Lifting.} Let $F=x^3+y^3+z^3-n$. A solution modulo $p$ has some coordinate nonzero modulo $p$, since $p\nmid n$; say $p\nmid x$, so $\partial F/\partial x=3x^2\not\equiv0\pmod p$ (here $p\ne3$). Let $(x,y,z)$
solve $F\equiv0\pmod{p^j}$ with $j\ge1$ and $p\nmid x$. The points congruent to it modulo $p^j$ are $(x+p^ja,y+p^jb,z+p^jc)$ with $a,b,c$ modulo $p$, and by Taylor's formula (all higher terms are
divisible by $p^{2j}$, hence by $p^{j+1}$)
\[
F(x+p^ja,y+p^jb,z+p^jc)\equiv F(x,y,z)+p^j\bigl(3x^2a+3y^2b+3z^2c\bigr)\pmod{p^{j+1}}.
\]
Dividing by $p^j$, the condition $F\equiv0\pmod{p^{j+1}}$ is one linear equation in $(a,b,c)$ over $\mathbb F_p$ with a nonzero coefficient $3x^2$, which has exactly $p^2$ solutions. So each solution modulo $p^j$ has exactly $p^2$
lifts modulo $p^{j+1}$, all with $p\nmid x$ again. By induction there are $p^2\cdot(p^2)^{k-1}=p^{2k}$ solutions modulo $p^k$, and the ratio in the definition of $\sigma_p(n)$ is $1$ for every $k$.

(b) By the proof of Proposition~\ref{prop:mod9}, there is no solution modulo $9$, hence none modulo $3^k$ for $k\ge2$.
\end{proof}

For the primes $p\equiv1\pmod3$, cubing is three-to-one on $\mathbb F_p^\times$ and $\sigma_p(n)$ is not $1$. For $p\nmid3n$ the lifting argument above still applies, so $\sigma_p(n)$ is the count modulo $p$ divided by $p^2$, and the theory of Jacobi sums shows that it is $1+O(1/p)$. Numerically, one finds a striking pattern: $p(\sigma_p(n)-1)$ is close to $6$ when $n$ is a cube modulo $p$ and close to $-3$ when it is not. A random $n$ is a cube modulo $p$ with probability $\frac13$, and $\frac13\cdot6+\frac23\cdot(-3)=0$: the corrections cancel on average, and the product over $p$ converges, though only conditionally and slowly. Computing it accurately is a small research problem in itself, which Booker and Sutherland solved.

\begin{remark}[Heath-Brown's conjecture, 1992]\label{rem:hb}
Combining the two factors, Heath-Brown conjectured that for $n\not\equiv\pm4\pmod9$ and $n$ not a cube, the number of solutions with $\max(|x|,|y|,|z|)\le B$, counted up to permutation, satisfies
\[
N_n(B)\sim\rho(n)\log B,\qquad\rho(n)=\frac16\,\sigma_\infty\prod_p\sigma_p(n)>0.
\]
In particular every such $n$ would be a sum of three cubes in infinitely many ways. This is the same shape of formula that the circle method \emph{proves} for Waring's problem with
enough variables (Tome~13, Theorem~7.11, with its singular series and singular integral); for three cubes the method is far out of reach, and the formula is only a prediction.
\end{remark}

The heuristic predicts more than existence. Since solutions are spread evenly over the scales $\log B$, the smallest solution for a given $n$ has $\log\max(|x|,|y|,|z|)$ of order $1/\rho(n)$,
with large fluctuations: some $n$ have a solution with three-digit coordinates, others need sixteen digits or more. The record sizes of the solutions for $30$, $33$ and $42$ are not
signs that something unusual happens for these numbers; they are what a constant density on a logarithmic scale produces, once in a while. Booker and Sutherland compared the counts
of all known solutions for $n\le1000$ with the predicted $\rho(n)\log B$ and found close agreement.

\section{How the search was done}

The naive search for solutions with coordinates up to $B$ looks at $B^2$ pairs $(x,y)$ and asks whether $n-x^3-y^3$ is a cube: hopeless for $B=10^{16}$. Booker's search is almost linear in $B$, and
it finds all solutions whose \emph{smallest} coordinate is at most $B$, however large the other two are. It rests on one factorisation.

\begin{proposition}[the divisor trick]\label{prop:booker}
Let $n$ be an integer that is not a cube, and let $x^3+y^3+z^3=n$. Put $d=x+y$. Then $d\ne0$, $d$ divides $n-z^3$,
\[
|d|^3\le4\,|n-z^3|,\qquad\text{and}\qquad\Delta=\frac{4(n-z^3)-d^3}{3d}=(x-y)^2.
\]
Conversely, if $d\ne0$ and $z$ are integers with $\Delta$ a perfect square, and $x,y=\frac{d\pm\sqrt\Delta}2$ are integers, then $x^3+y^3+z^3=n$.
\end{proposition}

\begin{proof}
If $d=0$ then $y=-x$ and $z^3=n$, which is excluded. From $x^3+y^3=(x+y)(x^2-xy+y^2)$ we get $n-z^3=d(x^2-xy+y^2)$, so $d\mid n-z^3$. Since $x^2-xy+y^2=\frac14\bigl((x+y)^2+3(x-y)^2\bigr)\ge\frac{d^2}4$, we get
$|n-z^3|\ge\frac{|d|^3}4$. Next, $x^3+y^3=d^3-3dxy$, so $xy=\frac{d^3-(n-z^3)}{3d}$, and $(x-y)^2=d^2-4xy=\frac{3d^3-4d^3+4(n-z^3)}{3d}=\Delta$. Conversely, if $x+y=d$ and $(x-y)^2=\Delta$, then $xy=\frac{d^2-\Delta}4=\frac{d^3-(n-z^3)}{3d}$,
and $x^3+y^3=d^3-3dxy=n-z^3$.
\end{proof}

\begin{corollary}\label{cor:booker}
Let $n$ be not a cube, and $B\ge1$. Every solution of $x^3+y^3+z^3=n$ with $\min(|x|,|y|,|z|)\le B$ is found by the following search: for each nonzero $d$ with $|d|\le(4|n|+4B^3)^{1/3}$, for each
$z$ with $|z|\le B$ and $z^3\equiv n\pmod{|d|}$, test whether $\Delta$ is a square. The number of pairs $(d,z)$ examined is $\sum_{d}r(d)\bigl(\frac{2B}{|d|}+1\bigr)$, where $r(d)$ is the number of
cube roots of $n$ modulo $|d|$.
\end{corollary}

\begin{proof}
Permute the coordinates so that $|z|$ is the smallest of $|x|,|y|,|z|$, so $|z|\le B$; then Proposition~\ref{prop:booker} gives $d\ne0$ with $d\mid n-z^3$, so $z^3\equiv n\pmod{|d|}$, and $|d|^3\le4|n-z^3|\le4|n|+4B^3$. The solution is then recovered from
$(d,z)$. The count is that of integers $z\in[-B,B]$ in $r(d)$ residue classes modulo $|d|$.
\end{proof}

Since $r(d)$ is small on average (it is $1$ when every prime factor of $d$ is $\equiv2\pmod3$ and prime to $3n$), the sum is about $B$ times a power of $\log B$: almost linear. Booker
added a sieve to discard most $z$ before computing $\Delta$, and with Sutherland a dozen further refinements, including a use of cubic reciprocity. In the solution for
$33$, the smallest coordinate is $z=-2\,736\,111\,468\,807\,040$, and $d=x+y=87\,723\,532\,425\,289$ is about thirty times smaller than $|z|$ (Exercise~\ref{ex:booker33}). So this solution sits at a value of $d$ far below the bound $(4|n|+4B^3)^{1/3}\approx1.6\,|z|$ allowed by the corollary.

\begin{think}
Why does it help to search by the \emph{smallest} coordinate? In the third representation of $3$, how large is the largest coordinate compared with the smallest?
\end{think}

\begin{beauty}
Three cubes, a mod $9$ obstruction, a logarithmic density, a factorisation of $x^3+y^3$ known to Euler, and half a million volunteer computers. The heuristic did not prove that $42$ is a sum of
three cubes; it said that the solutions were probably there, how large they would probably be, and how to search for them, and it was right on all three counts. Mathematics has few
predictions so precise about objects so stubborn.
\end{beauty}

\begin{remark}[open cases]\label{rem:open3cubes}
After the work of Booker and Sutherland, the integers below $1000$ not $\equiv\pm4\pmod9$ for which no representation was known were $114$, $390$, $627$, $633$, $732$, $921$ and $975$. The heuristic predicts
solutions for all of them. That every $n\not\equiv\pm4\pmod9$ is a sum of three cubes is open; so is whether $3$ has infinitely many representations. For the deeper question of why the
heuristic should hold, the relevant concept is the Brauer--Manin obstruction for integral points on cubic surfaces; J.-L.~Colliot-Th\'el\`ene and O.~Wittenberg showed in 2012 that it does not
obstruct sums of three cubes.
\end{remark}

\section{Exercises}

\begin{exercise}[$\star$]\label{ex:mod7}
Show that the cubes modulo $7$ are $0,\pm1$, as modulo $9$, but that every residue modulo $7$ is a sum of three cubes. Which moduli $m$ give an obstruction?
\end{exercise}

\begin{exercise}[$\star$]\label{ex:cubebij}
Show that $x\mapsto x^3$ is a bijection of $\Z/2\Z$ but not of $\Z/8\Z$. Why does the proof of Proposition~\ref{prop:localdens}(a) lift solutions from $\mathbb F_p$ instead of using a bijection modulo $p^k$?
\end{exercise}

\begin{exercise}[$\star\star$]\label{ex:dimcount}
\emph{Counting dimensions.} Repeat the real-density heuristic for $x_1^3+\dots+x_k^3=n$ and show that it predicts about $B^{k-3}$ solutions up to $B$ (and $\log B$ for $k=3$). Then show, without
any heuristic, that $x^3+y^3=n$ has only finitely many integer solutions for $n\ne0$.
\end{exercise}

\begin{exercise}[$\star$]\label{ex:families2}
Use Proposition~\ref{prop:families} to show that $2m^3$ and $m^3$ have infinitely many representations as sums of three cubes for every $m\ne0$. Is $16$ a sum of three cubes in infinitely many ways?
\end{exercise}

\begin{exercise}[$\star\star$]\label{ex:booker33}
For the representation of $33$ above, compute $d=x+y$ (for the two largest coordinates), check that $d\mid33-z^3$, compute $\Delta$ and recover $x$ and $y$ from $(d,z)$ as in Proposition~\ref{prop:booker}. How does $|d|$ compare
with the bound $(4|33-z^3|)^{1/3}$?
\end{exercise}

\begin{exercise}[$\star\star\star$]\label{ex:sigmainf}
\emph{The real density.} For $t>0$, the number of integer $z$ with $|z^3-t^3|\le\frac12$ is about the length of the interval of such real $z$. Show that this length is $\frac1{3t^2}(1+o(1))$ as $t\to\infty$.
Deduce that the expected number of solutions of $x^3+y^3+z^3=n$ with $B<|z|\le2B$ and $|x|,|y|\le|z|$, summed over integer $(x,y)$, is about
\[
\iint\frac{\dd x\dd y}{3|x^3+y^3|^{2/3}}
\]
over the corresponding region, and that this is independent of $B$ by scaling. (Computing the constant exactly gives $\sigma_\infty$; you need not do it.)
\end{exercise}

\section{Hints}

\begin{hint}{ex:mod7}
Sums of three elements of $\{0,\pm1\}$ cover $\{-3,\dots,3\}$. For an obstruction modulo $m$ the cubes must be rare: try $m=9$, then use Cauchy--Davenport and lifting for primes $p\ne3$.
\end{hint}

\begin{hint}{ex:cubebij}
$2^3\equiv0\pmod8$.
\end{hint}

\begin{hint}{ex:dimcount}
Scale $x_i=Bu_i$; the volume of $\{|F-n|<\frac12\}$ in a shell scales like $B^{k}\cdot B^{-3}$. For two cubes, $d=x+y$ divides $n$.
\end{hint}

\begin{hint}{ex:families2}
Multiply the identities by $m^3$.
\end{hint}

\begin{hint}{ex:booker33}
Use a computer or exact arithmetic; $\Delta$ has $32$ digits.
\end{hint}

\begin{hint}{ex:sigmainf}
$(t^3+\frac12)^{1/3}-(t^3-\frac12)^{1/3}\approx\frac{\dd}{\dd s}s^{1/3}\big|_{s=t^3}$.
\end{hint}

\section{Solutions}

\begin{solution}{ex:mod7}
\idea Cubes modulo $m$ are rare when the cubing map is far from injective.
\writeup Modulo $7$: $1,2,4\mapsto1$ and $3,5,6\mapsto-1$ (as $2^3=8$, $3^3=27\equiv-1$, and so on), so the cubes are $0,\pm1$. Sums of three of them give $-3,\dots,3$, that is all residues modulo $7$. Modulo $9$ the same
three cubes cover only seven of the nine residues, which is why $\pm4$ are excluded.

Now let $p\ne3$ be prime. Modulo $p$ every residue is a sum of three cubes: for $p\equiv2\pmod3$ (or $p=2$) every residue is a cube, and for $p\equiv1\pmod3$ the set $C$ of cubes has $\frac{p-1}3+1$ elements, so the Cauchy--Davenport theorem (Tome~2, Theorem~M.11) gives $|C+C+C|\ge\min(p,3|C|-2)=p$. Each residue has a solution with a coordinate prime to $p$ (for $0$ take $1+(-1)+0$), and the lifting step of Proposition~\ref{prop:localdens} carries it to every $p^k$. By the Chinese remainder theorem, the only possible obstructions live at powers of $3$; a direct computation verifies that modulo $27$ and $81$ the excluded residues are exactly those $\equiv\pm4\pmod9$. (Lifting further needs the stronger Hensel lemma for a derivative divisible by $3$ exactly once; try it.) So $9$ is the modulus that matters.
\end{solution}

\begin{solution}{ex:cubebij}
\idea A bijection of a finite set is the same as an injection.
\writeup Modulo $2$, $0^3=0$ and $1^3=1$. Modulo $8$, $0^3\equiv2^3\equiv4^3\equiv6^3\equiv0$, so cubing is not injective. In the proof, only the solutions modulo $p$ were counted by a bijection; the lifting to $p^k$ used the
nonvanishing of a partial derivative, which holds because $p\nmid n$ forces a coordinate to be a unit. A bijection argument modulo $p^k$ would be false.
\end{solution}

\begin{solution}{ex:dimcount}
\idea Each scale contributes $B^{k}$ points, of which a proportion $B^{-3}$ hit a given value.
\writeup In the shell $B<\max|x_i|\le2B$ there are about $B^k$ integer points, and $F=\sum x_i^3$ takes values of size $B^3$ there, spread over an interval of length about $B^3$; the expected number with $F=n$ is
about $B^{k-3}$ times a constant (the scaled density, as in the text). Summing over dyadic shells up to $B$ gives $B^{k-3}$ for $k\ge4$, $\log B$ for $k=3$, and a convergent sum for $k\le2$: finitely many.
For $k=2$ this is a theorem: if $x^3+y^3=n\ne0$, then $d=x+y$ is a nonzero divisor of $n$, and $x^2-xy+y^2=\frac nd$; and $\bigl|\frac nd\bigr|\le|n|$. As $x^2-xy+y^2=\frac12(x^2+y^2)+\frac12(x-y)^2\ge\frac12(x^2+y^2)$, we get $x^2+y^2\le2|n|$. So there are finitely many solutions.
\end{solution}

\begin{solution}{ex:families2}
\idea Cubes are homogeneous.
\writeup Multiplying $(1+6t^3)^3+(1-6t^3)^3+(-6t^2)^3=2$ by $m^3$ gives $(m+6mt^3)^3+(m-6mt^3)^3+(-6mt^2)^3=2m^3$ for every $t$, and different $t$ give different triples. Similarly Mahler's identity gives infinitely
many representations of $m^3$. So $16=2\cdot2^3$ has infinitely many, for example $t=1$ gives $14^3+(-10)^3+(-12)^3=2744-1000-1728=16$.
\end{solution}

\begin{solution}{ex:booker33}
\idea Plug into the formulas.
\writeup With $x=8\,866\,128\,975\,287\,528$, $y=-8\,778\,405\,442\,862\,239$, $z=-2\,736\,111\,468\,807\,040$: $d=x+y=87\,723\,532\,425\,289$. Exact arithmetic confirms $d\mid33-z^3$, that
$\Delta=(x-y)^2$ with $x-y=17\,644\,534\,418\,149\,767$, and that $\frac{d\pm\sqrt\Delta}2$ gives back $x$ and $y$. The bound $(4|33-z^3|)^{1/3}=4^{1/3}|z|(1+o(1))\approx4.34\times10^{15}$ is about fifty times larger than $d$: the actual
$d$ is much smaller than the worst case allowed by the proposition, which is what makes the refined search of Booker and Sutherland pay off.
\end{solution}

\begin{solution}{ex:sigmainf}
\idea Linearise the cube root.
\writeup By the mean value theorem, $(t^3+\frac12)^{1/3}-(t^3-\frac12)^{1/3}=\frac13s^{-2/3}$ for some $s$ with $|s-t^3|\le\frac12$, and $s^{-2/3}=t^{-2}(1+O(t^{-3}))$. So the length is $\frac1{3t^2}(1+o(1))$. For fixed $(x,y)$ the
relevant $t$ is $|n-x^3-y^3|^{1/3}$, which is $|x^3+y^3|^{1/3}(1+o(1))$ when $|x^3+y^3|$ is large compared with $n$; summing $\frac1{3t^2}$ over integer $(x,y)$ in the region and replacing the sum by an integral gives the
displayed integral. Substituting $x=Bu$, $y=Bv$ multiplies $\dd x\dd y$ by $B^2$ and $|x^3+y^3|^{2/3}$ by $B^2$, and maps the region for $B$ to the region for $1$; so the integral does not depend on $B$. (The integrand
is singular on the line $u=-v$, but like $|u+v|^{-2/3}$, which is integrable.)
\end{solution}

\section*{Sources and further reading}
\addcontentsline{toc}{section}{Sources and further reading}

\begin{itemize}
\item \emph{The search.} A.~R.~Booker, ``Cracking the problem with 33'', \emph{Res.\ Number Theory} 5 (2019), paper 26. A.~R.~Booker and A.~V.~Sutherland, ``On a question of Mordell'',
\emph{Proc.\ Natl.\ Acad.\ Sci.\ USA} 118 (2021), e2022377118, arXiv:2007.01209 (the solutions for $3$ and $42$, Mordell's 1953 question, the density computations and the list of open cases).
S.~G.~Huisman, ``Newer sums of three cubes'', arXiv:1604.07746, 2016. ``Sum-of-three-cubes problem solved for `stubborn' number 33'', \emph{Quanta Magazine}, 26 March 2019, and the University of Bristol
press release of 6 September 2019, for the story and the Numberphile video.
\item \emph{The heuristic.} D.~R.~Heath-Brown, ``The density of zeros of forms for which weak approximation fails'', \emph{Math.\ Comp.}\ 59 (1992), 613--623.
\item \emph{Families and obstructions.} K.~Mahler, ``Note on hypothesis K of Hardy and Littlewood'', \emph{J.\ London Math.\ Soc.}\ 11 (1936), 136--138. J.-L.~Colliot-Th\'el\`ene and O.~Wittenberg,
``Groupe de Brauer et points entiers de deux familles de surfaces cubiques affines'', \emph{Amer.\ J.\ Math.}\ 134 (2012), 1303--1327.
\end{itemize}
